\label{sec:course}
Table~\ref{tab:course} maps each topic of the CSE 402 lab outline to the method we implemented
for it and to where that method is used. The project covers all four lab assignments and the
optimization home assignment, and so all three course outcomes: nonlinear equations and linear
systems (CO1), interpolation, regression, and numerical integration (CO2), and simulation with
random number generation and Monte Carlo methods (CO3). Unless marked $^\dagger$, every method in
the table is our own code in \texttt{core/sirlab}, which never imports SciPy (enforced by
\texttt{tests/test\_hygiene.py}); SciPy appears only in tests, as an independent cross-check.

\begin{table*}[tp]
  \centering
  \footnotesize
  \caption{CSE 402 syllabus topics and where the project uses them. Code paths are relative to
    \texttt{core/sirlab/}; \texttt{E}-numbered scripts are in \texttt{experiments/}.
    $^\dagger$\,Computed with a NumPy routine rather than our own code.}
  \label{tab:course}
  \begin{tabularx}{\textwidth}{@{}>{\raggedright\arraybackslash}p{2.7cm}
      >{\raggedright\arraybackslash}X >{\raggedright\arraybackslash}p{3.5cm}@{}}
    \toprule
    Syllabus topic & What we implemented, and where it is used & Code \\
    \midrule
    \multicolumn{3}{@{}l}{\textbf{Lab 1: approximations, errors, and root finding (CO1)}} \\
    Truncation error & Global error vs.\ the gold standard; observed order $\hat p$ (E01); bias in
      $\hat\beta \propto h^p$ (E07) & \texttt{diagnostics.py} \\
    Round-off error & $S+I+R=N$ holds to round-off for every explicit RK method at any $h$ (E02);
      adaptive Simpson stops refining at the round-off floor
      & \texttt{diagnostics.py}, \texttt{quadrature/adaptive.py} \\
    Bisection, Newton--Raphson & Safeguarded Newton (\texttt{rtsafe}): Newton steps inside the
      bracket, bisection otherwise; inverts $t(R)$ and solves the final-size equation
      & \texttt{roots/newton.py} \\
    Newton--Raphson for systems & $(I - hJ_f)\,\delta = -G(y)$ on each backward-Euler or
      trapezoidal stage & \texttt{solvers/implicit.py} \\
    Plotting & All figures generated from the experiment JSON artifacts
      & \texttt{scripts/make\_figures.py} \\
    \midrule
    \multicolumn{3}{@{}l}{\textbf{Lab 2: systems of linear equations and eigenvalues (CO1)}} \\
    Gauss elimination, LU decomposition & Doolittle LU with partial pivoting, forward/back
      substitution; every implicit Newton step, the GN/LM normal equations, every inverse-power step
      & \texttt{linalg/lu.py} \\
    Matrix inverse, conditioning & $(J^\top J)^{-1}$ for the asymptotic covariance via LU solves;
      1-norm $\kappa(J^\top J)$ (E05); 2-norm $\kappa(J)$ from the eigenvalues of $J^\top J$ (E10)
      & \texttt{linalg/lu.py}, \texttt{linalg/eigen.py} \\
    Power method & Largest-entry normalization; $h\,\rho(J_y f)$ for stability (E03), with a 2-D
      projection when the SIR Jacobian's eigenvalues are complex (Section~\ref{sec:eigen})
      & \texttt{linalg/eigen.py}, \texttt{diagnostics.py} \\
    Inverse power method & Smallest eigenvalue by LU solves, factored once; $\lambda_{\min}(J^\top J)$
      for $\kappa_2(J)$ (E10) & \texttt{linalg/eigen.py} \\
    Hotelling deflation & All eigenpairs of a symmetric matrix: $\mathrm{Cov}(\hat\theta)$ for the
      ellipse axes (E11, E14); $J^\top J$ for the cost valley (E05)
      & \texttt{linalg/eigen.py}, \texttt{uq/asymptotic.py} \\
    \midrule
    \multicolumn{3}{@{}l}{\textbf{Lab 3: random number generation and Monte Carlo methods (CO3)}} \\
    Random number generation & PCG64 uniform stream; exact normal, uniform, integer, and gamma
      variates; per-replicate seeds (Section~\ref{sec:rng})$^\dagger$
      & \texttt{observe.py}, E05--E11, E13 \\
    Monte Carlo simulation & Replicated synthetic datasets from the known truth, each refitted:
      bias, variance, RMSE (E07--E10, E13); UQ route 3 (E11)
      & \texttt{uq/montecarlo.py}, E07--E11, E13 \\
    Bootstrap resampling & Residual bootstrap of $(\hat\beta, \hat\gamma)$ (E11); percentile band
      on the RMSE curve (E08) & \texttt{uq/bootstrap.py}, \texttt{E08\_noise.py} \\
    Metropolis--Hastings & Adaptive random-walk Metropolis in $\ln\theta$ plus a Gibbs step for
      $\sigma^2$; 4 chains, acceptance rate, $\hat R$~\cite{gelman1992}, ESS (E11)
      & \texttt{uq/mcmc.py} \\
    \midrule
    \multicolumn{3}{@{}l}{\textbf{Home Assignment 1: QR and optimization (CO1)}} \\
    QR & Householder QR step $R\,\delta = -Q^\top r$ avoids the $\kappa(J)^2$ of the normal
      equations; unit-tested, but reported fits use LU. Not the QR eigenvalue iteration
      & \texttt{linalg/qr.py}, \texttt{estimate/gaussnewton.py} \\
    Golden-section search & 1-D minimizer in coordinate descent (E06) and the profile likelihoods
      (E11) & \texttt{estimate/golden.py}, \texttt{uq/profile.py} \\
    Newton's method & Gauss--Newton: Newton with Hessian $\approx J^\top J$, backtracking line
      search (E06) & \texttt{estimate/gaussnewton.py} \\
    Gradient methods & Levenberg--Marquardt spans Gauss--Newton ($\lambda \to 0$) and scaled
      gradient descent ($\lambda \to \infty$); the estimator in E07--E11, E13, E14
      & \texttt{estimate/lm.py} \\
    Constrained optimization & $\beta, \gamma > 0$ by optimizing $u = \ln\theta$; box bounds in
      golden-section search & \texttt{estimate/lm.py}, \texttt{estimate/golden.py} \\
    Beyond the list & Nelder--Mead: derivative-free baseline (E06) and 3-parameter fit (E10);
      $J(\beta, \gamma)$ on a 2-D grid (E05)
      & \texttt{estimate/neldermead.py}, \texttt{E05\_landscape.py} \\
    \midrule
    \multicolumn{3}{@{}l}{\textbf{Lab 4: curve fitting, interpolation, and numerical integration (CO2)}} \\
    Least-squares regression & Nonlinear least squares \eqref{eq:cost} in every fit; linear least
      squares on $(\ln h, \ln e)$ gives $\hat p$ (E01) & \texttt{estimate/}, \texttt{diagnostics.py} \\
    Linear interpolation & Locates the crossover $\sigma^*$ between grid points (E13)
      & \texttt{E13\_crossover.py} \\
    Spline interpolation & Piecewise cubic Hermite dense output evaluates trajectories and
      sensitivities at the observation times & \texttt{solvers/base.py} \\
    Trapezoidal and Simpson's rules & Adaptive Simpson for the gold-standard integral
      \eqref{eq:tofr}. Composite and uneven-grid Simpson (the integral of the 3-point Lagrange
      quadratic) and the trapezoidal rule serve the incidence operator: unit-tested, but unused in
      the reported experiments & \texttt{quadrature/adaptive.py}, \texttt{quadrature/simpson.py} \\
    Richardson, Romberg & Each accepted Simpson panel adds $(S_2 - S_1)/15$ (one- vs.\ two-panel
      estimates): one Romberg step & \texttt{quadrature/adaptive.py} \\
    Numerical differentiation & Central-difference node slopes for the Hermite
      interpolant$^\dagger$; complex-step checks of the analytic Jacobians
      & \texttt{solvers/base.py}, \texttt{tests/test\_models.py} \\
    \midrule
    \multicolumn{3}{@{}l}{\textbf{Beyond the lab outline: ODE solvers (course textbook~\cite{chapra2015}, Part Seven)}} \\
    Runge--Kutta methods & Euler, Heun, and classical RK4, with their stability frontier (E03);
      adaptive Dormand--Prince RK45 with PI step control
      & \texttt{solvers/explicit.py}, \texttt{solvers/rk45.py} \\
    Implicit methods & Backward Euler, trapezoidal rule (A-stable), in the convergence study (E01)
      & \texttt{solvers/implicit.py} \\
    \bottomrule
  \end{tabularx}
\end{table*}

Two further points complete the picture. First, the ODE solvers are not a CSE 402 lab topic: we
took them from Part Seven of the course textbook~\cite{chapra2015}, and they are the project's
experimental variable rather than its apparatus. Second, a few syllabus methods have no role in
this problem and are not used: false position, Bairstow's method, Gauss--Jordan elimination, the
QR eigenvalue iteration, and Newton's divided-difference interpolation.
